\documentclass[12pt,a4paper]{article}
\usepackage[utf8]{inputenc}
%Set page layout parameters
\usepackage[
top=4cm,
bottom=2cm,
left=2cm,
right=2cm,
headheight=2cm,
footskip=1cm,
marginparwidth=2cm]{geometry}
\usepackage{amsmath}
\usepackage[dvipsnames]{xcolor}
% To use color boxes
\usepackage{tcolorbox}
\usepackage{lipsum}
\tcbuselibrary{listings,theorems,skins,vignette,raster,breakable,magazine,poster,fitting,hooks}
\newcounter{ejemplo}
%Images path
\graphicspath{{images/}}
\usepackage{float} %provide us with the option H
\usepackage{fancyhdr}
\usepackage{tikz}
\usepackage{tikzpagenodes}
%to redefine figure caption 
\renewcommand{\figurename}{Figura}
%change text font
\usepackage[T1]{fontenc}
\usepackage{tgbonum} % LaTeX will use the TEX Gyre Bonum font family

\definecolor{subtitle_yellow}{RGB}{252,213,129}
\definecolor{title_yellow}{RGB}{249,175,16}
\definecolor{steel_blue_title}{RGB}{93,138,168}
\definecolor{pewter_blue_subtitle}{RGB}{149,178,198}
\definecolor{ecru_circle}{RGB}{215,190,130}
\definecolor{ecru_circle1}{RGB}{68,148,173}

\newcommand{\myfirststyle}{%
	\begin{tikzpicture}[remember picture, overlay]
		
		%top border
		\fill[fill=steel_blue_title](current page.north west) rectangle ([yshift=-15mm]current page.north east);
		
		\fill[fill=pewter_blue_subtitle]([yshift=-17mm]current page.north west) rectangle ([yshift=-30mm]current page.north east);
		
		%page box
		%\draw[lightgray, fill=Rhodamine, fill opacity=0.05, line width=1mm, rounded corners=3mm] ([xshift=-5mm,yshift=5mm]current page text area.north west) rectangle ([xshift=5mm,yshift=-5mm]current page text area.south east);
		
		%bottom box
		\fill[fill=pewter_blue_subtitle](current page.south west) rectangle ([yshift=10mm]current page.south east);
		
		%page circle
		\draw[white,fill=ecru_circle1,line width=2mm] ([xshift=-5mm,yshift=25mm]current page text area.north east) circle (1); 
		
		
		%page number text        
		%\node[anchor=center] at ([xshift=-5mm,yshift=25mm]current page text area.north east) {\Huge\bfseries\sffamily\color{white}{\thepage}};
		
		%title text
		\node[anchor=west] at ([yshift=32mm]current page text area.north west) {\Huge\bfseries\sffamily\color{white}{Amplificador Inversor}};
		
		%subtitle text
		\node[anchor=west] at ([xshift=10mm,yshift=16mm]current page text area.north west) {\Large\bfseries\sffamily\color{white}{}};
		
		%footer text
		\node[anchor=center] at ([yshift=-15mm]current page text area.south) {\large\bfseries\sffamily\color{white}{Created by: Rodrigo Castillejos}};
		
	\end{tikzpicture}
}


%create page style
\fancypagestyle{MyFirstStyle}{%
	\fancyhf{}
	\fancyhead[C]{\myfirststyle}
	\renewcommand{\headrulewidth}{0pt}
}%

\pagestyle{MyFirstStyle}

\colorlet{colexam}{red!75!black}
\newtcolorbox[use counter=ejemplo]{myexample}{%
	empty,title={Ejemplo \thetcbcounter},attach boxed title to top left,
	boxed title style={empty,size=minimal,toprule=2pt,top=4pt,
		overlay={\draw[colexam,line width=2pt]
			([yshift=-1pt]frame.north west)--([yshift=-1pt]frame.north east);}},
	coltitle=colexam,fonttitle=\Large\bfseries,
	before=\par\medskip\noindent,parbox=false,boxsep=0pt,left=0pt,right=3mm,top=4pt,
	breakable,pad at break*=0mm,vfill before first,
	overlay unbroken={\draw[colexam,line width=1pt]
		([yshift=-1pt]title.north east)--([xshift=-0.5pt,yshift=-1pt]title.north-|frame.east)
		--([xshift=-0.5pt]frame.south east)--(frame.south west); },
	overlay first={\draw[colexam,line width=1pt]
		([yshift=-1pt]title.north east)--([xshift=-0.5pt,yshift=-1pt]title.north-|frame.east)
		--([xshift=-0.5pt]frame.south east); },
	overlay middle={\draw[colexam,line width=1pt] ([xshift=-0.5pt]frame.north east)
		--([xshift=-0.5pt]frame.south east); },
	overlay last={\draw[colexam,line width=1pt] ([xshift=-0.5pt]frame.north east)
		--([xshift=-0.5pt]frame.south east)--(frame.south west);},%
}


\begin{document}
	\section{Filtro pasa bajas de segundo orden}
	\subsection{Preliminares}
	La función de transferencia de un filtro pasa bajas de segundo orden está dado por
	\begin{tcolorbox}[colback=red!5!white,colframe=red!75!black,width=(\linewidth+0.45cm)]
		\begin{equation}
			H(s) = \frac{\omega_0^{2}}{s^2 + \frac{\omega_0}{Q} + \omega_0^{2} }
		\end{equation}
	\end{tcolorbox}
	Aquí $\omega_0$ es la frecuencia de esquina (o frecuencia central o frecuencia resonante) en radianes por segundo, y $Q$ es el factor de calidad.\\\\
	La ubicación de los polos son las raíces de $s^2 + \frac{\omega_0}{Q} + \omega_0^{2} = 0$
	\begin{align}
		p_{1,2} = -\frac{\omega_0}{2Q}\pm\frac{1}{2}\sqrt{\frac{\omega_0^{2}}{Q^2} - 4\omega_0^{2} } = -\frac{\omega_0}{2Q}\pm\frac{1}{2}\sqrt{ \omega_0^{2}\left( \frac{1}{Q^2} - 4 \right) }\nonumber
	\end{align}
	\begin{tcolorbox}[colback=red!5!white,colframe=red!75!black,width=(\linewidth+0.45cm)]
		\begin{subequations}
			\begin{align}
				p_1 &= - \frac{\omega_0}{2Q} + \frac{\omega_0}{2Q}\sqrt{1-4Q^2}\\
				p_2 &= - \frac{\omega_0}{2Q} + \frac{\omega_0}{2Q}\sqrt{1-4Q^2}
			\end{align}	
		\end{subequations}
	\end{tcolorbox}
	Para que los polos sean complejos conjugados el resultado dentro de la raíz debe ser negativo, para esto, resolvemos la ecuación que se ubica dentro de la raíz
	
	\begin{align}
		1-4Q^2&<0\nonumber\\
		1&<4Q^2\nonumber
	\end{align}
	Es decir
	\begin{tcolorbox}[colback=red!5!white,colframe=red!75!black,width=(\linewidth+0.45cm)]
		\begin{equation}
			Q > \frac{1}{2}
		\end{equation}
	\end{tcolorbox}
	Usualmente es de principal interés el caso de los polos complejos conjugados, obtenidos por $Q>0.5$. La figura 1 muestra la ubicación del par de polos complejos conjugados en el plano $s$. Observe que la distancia radial de los modos naturales (desde el origen) es igual a $\omega_0$ la cual es conocida como la frecuencia del polo.\\
	El parámetro $Q$ determina la distancia de los polos desde el eje $j\omega$. Cuanto mayor sea el valor de $Q$, más cerca estarán los polos del eje $j\omega$ y la respuesta del filtro se vuelve más selectiva. Un valor infinito para $Q$ ubica a los polos en el eje $j\omega$ y puede producir oscilaciones sostenidas en la realización del circuito. Un valor negativo de $Q$ implica que los polos están en la mitad derecha del plano $s$, lo que ciertamente produce oscilaciones.
	\begin{figure}[H]
		\centering
		\includegraphics[scale=1]{figure1}
		\caption{Definición de los parámetros $\omega_0$ y $Q$ de un par de polos complejos conjugados.}
		\label{fig:pic1}
	\end{figure}
	
	Los ceros de un filtro de segundo orden están determinados por los coeficientes del numerador $a_0, a_1$ y $a_2$. Así, los coeficientes del numerador determinan el tipo de filtro de segundo orden (es decir, LP,HP, etc.)
	
	Cuando los polos son complejos conjugados, tenemos
	\begin{tcolorbox}[colback=red!5!white,colframe=red!75!black,width=(\linewidth+0.45cm)]
		\begin{subequations}
			\begin{align}
				p_1 &= - \frac{\omega_0}{2Q} + j\omega_0\sqrt{1-4Q^2}\\
				p_2 &= - \frac{\omega_0}{2Q} - j\omega_0\sqrt{1-4Q^2}
			\end{align}	
		\end{subequations}
	\end{tcolorbox}

	Conforme $Q$ disminuye a cero $(Q\rightarrow 0)$, de la ecuación (2a) y (2b) tenemos que $p_1 \rightarrow 0$ y $p_2 \rightarrow - \infty$
	
	\subsection{Respuesta en frecuencia}
	La respuesta en frecuencia del filtro pasa bajas de segundo orden se obtiene sustituyendo $s=j\omega$ en $H(s)$.
	
	\begin{align}
		H(\omega) = H(s)|_{s=j\omega} = \frac{\omega_0^{2}}{(j\omega)^{2} + \frac{\omega_0}{Q}j\omega + \omega_0^{2}}\nonumber
	\end{align}
	\begin{tcolorbox}[colback=red!5!white,colframe=red!75!black,width=(\linewidth+0.45cm)]
		\begin{align}
			H(\omega) = \frac{\omega_0^{2}}{-\omega^{2} + \frac{\omega_0}{Q} + \omega_0^{2}}
		\end{align}	
	\end{tcolorbox}

	\subsection{Respuesta en magnitud}
	La respuesta en magnitud del filtro pasa bajas de segundo orden está dado por el valor absoluto de $H(\omega)$.
	\begin{align}
		\left| H(\omega) \right| = \frac{\omega_0^{2}}{\left|\omega^{2}  + \frac{\omega_0}{Q}j\omega + \omega_0^{2} \right|}\nonumber
	\end{align}
	\begin{tcolorbox}[colback=red!5!white,colframe=red!75!black,width=(\linewidth+0.45cm)]
		\begin{align}
			|H(\omega)| = \frac{\omega_0^{2}}{\sqrt{(\omega_0^{2} - \omega^{2})^{2} + (\frac{\omega_0}{Q}\omega)^{2} }}
		\end{align}	
	\end{tcolorbox}
	
	En $\omega = \omega_0$ la magnitud es
	\begin{align}
		|H(\omega)| = \frac{ \omega_0^{2} }{ \sqrt{ (\omega_0^{2} - \omega_0^{2} ) + \left(  \frac{\omega_0^{2}}{Q}\right)^{2} } } = \frac{\omega_0^{2}}{\frac{\omega_0^{2}}{Q}} = Q\nonumber 
	\end{align}

	La frecuencia pico $\omega_{pico}$ donde la magnitud es máxima, está dada por
	\begin{align}
		\frac{d}{d\omega}|H(\omega)| = \frac{d}{d\omega} \frac{\omega_0^{2}}{\sqrt{(\omega_0^{2} - \omega^{2})^{2} + (\frac{\omega_0}{Q}\omega)^{2} }} = 0\nonumber\\
		\frac{d}{d\omega}|H(\omega)|^{2} = \frac{d}{d\omega}\frac{\omega_0^{4}}{(\omega_0^{2} - \omega^{2})^{2} + (\frac{\omega_0}{Q}\omega)^{2}} = 0\nonumber\\
		\frac{d}{d\omega}|H(\omega)|^{2} = \frac{d}{d\omega}\omega_0^{4}\left[ \left(\omega_0^{2} - \omega^{2}\right)^{2} + \left(\frac{\omega_0}{Q}\omega\right)^{2} \right]^{-1}\nonumber\\
		-\omega_0^{4}\left[ (\omega_0^{2} - \omega^{2}) + \left( \frac{\omega_0}{Q}\omega \right)^{2} \right]^{-2} \left[ 2(\omega_0^{2} - \omega^{2})(-2\omega) + 2\omega\frac{\omega_0^{2}}{Q^2} \right] = 0\nonumber	
	\end{align}
	
	\vspace{-0.5cm}

	\begin{align}
		-4\omega(\omega_0^{2} - \omega^{2}) + 2\omega\frac{\omega_0^{2}}{Q^2} = 0\nonumber\\
		-4\omega\omega_0^{2} + 4\omega^{3} + 2\omega\frac{\omega_0^{2}}{Q^2} = 0\nonumber\\
		\omega^{2} - \omega_0^{2} + \frac{\omega_0^{2}}{2Q^2} = 0\nonumber\\
		\omega^{2} = \omega_0^{2}\left( 1 - \frac{1}{2Q^{2}} \right)\nonumber
	\end{align}

	\begin{tcolorbox}[colback=red!5!white,colframe=red!75!black,width=(\linewidth+0.45cm)]
		\begin{align}
			\omega_{pico} = \omega_0 \sqrt{1 - \frac{1}{2Q^2}}
		\end{align}	
	\end{tcolorbox}
	
	La magnitud máxima en la frecuencia pico se obtiene sustituyendo la ecuación (7) en la ecuación (6)
		
	\begin{align}
		|H(j\omega_{pico})| &= \frac{\omega_0^2}{\sqrt{\left[ \omega_0^2  - \omega_0^2 \left( 1 - \frac{1}{2Q^2} \right) \right]^2 + \frac{\omega_0^2}{Q^2}\omega_0^2\left( 1 - \frac{1}{2Q^2} \right) } }\nonumber\\
		&=\frac{\omega_0^2}{ \sqrt{\left[ \omega_0^2 - \omega_0^2 + \frac{\omega_0^2}{2Q^2} \right]^2 + \frac{\omega_0^4}{Q^2} - \frac{\omega_0^4}{2Q^4} } }\nonumber
	\end{align}
	\begin{align}
		|H(j\omega_{pico})| = \frac{\omega_0^2}{ \sqrt{ \frac{\omega_0^4}{4Q^4} + \frac{\omega_0^4}{Q^2} - \frac{\omega_0^4}{2Q^4} } } = \frac{\omega_0^4}{ \sqrt{ \frac{\omega_0^4 + 4Q^2\omega_0^4 - 2\omega_0^4}{4Q^4} } } = \frac{\omega_0^4}{ \sqrt{ \frac{\omega_0^4(4Q^2 - 1)}{4Q^4} } } = \frac{\omega_0^2}{ \frac{\omega_0^2\sqrt{4Q^2 - 1}}{2Q^2} }\nonumber
	\end{align}
	\begin{align}
		|H(j\omega_{pico})| = \frac{2Q^2}{ \sqrt{4Q^2 - 1} } = \frac{2Q^2}{\sqrt{4Q^2\left( 1 - \frac{1}{4Q^2} \right) }} = \frac{2Q^2}{2Q\sqrt{1 - \frac{1}{4Q^2}} }\nonumber
	\end{align}

	Por tanto
	\begin{tcolorbox}[colback=red!5!white,colframe=red!75!black,width=(\linewidth+0.45cm)]
		\begin{align}
			|H(j\omega_{pico})| = K\frac{Q}{\sqrt{1 - \frac{1}{4Q^2}}}
		\end{align}	
	\end{tcolorbox}
	
	La frecuencia de corte se obtiene igualando la ecuación (6) con $1/\sqrt{2}$
	
	\begin{align}
		\frac{\omega_0^2}{ \sqrt{(\omega_0^2 - \omega^2)^2 + \left( \frac{\omega_0}{Q} \right)^2} } = \frac{1}{\sqrt{2}}
	\end{align}
	Elevando ambos lados de la ecuación (9) al cuadrado obtenemos
	\begin{align}
		\frac{\omega_0^4}{(\omega_0^2 - \omega^2)^2 + \left( \frac{\omega_0}{Q} \right)^2} = \frac{1}{2}\nonumber\\
		2\omega_0^4 = (\omega_0^2 - \omega^2)^2 + \left( \frac{\omega_0}{Q} \right)^2\nonumber\\
		2\omega_0^4 = \omega_0^4 - 2\omega^2\omega_0^2 + \omega^4 + \frac{\omega^2\omega_0^2}{Q^2}\nonumber\\
		\omega_0^4 + \frac{\omega_0^4}{Q^2}\omega^2 - 2\omega_0^2\omega^2 - \omega_0^4 = 0\nonumber\\
		\omega^4 - \left( 2\omega_0^2 - \frac{\omega_0^2}{Q^2} \right)\omega^2 - \omega_0^4 = 0\nonumber
	\end{align}

	Sea $x = \omega^2$ y $x^2 = \omega^4$, así
	\begin{align}
		x^2 - \left( 2\omega_0^2 - \frac{\omega_0^2}{Q^2} \right)x - \omega_0^4 = 0\nonumber\\
		x_{1,2} = \frac{\left( 2\omega_0^2 - \frac{\omega_0^2}{Q^2} \right) \pm \sqrt{\omega_0^4\left( 2 - \frac{1}{Q^2} \right)^2 + 4\omega_0^4 } }{2}\nonumber\\
		x_{1,2} = \frac{\omega_0^2\left( 2 - \frac{1}{Q^2} \right)\pm \omega_0^2\sqrt{\left( 2 - \frac{1}{Q^2} \right)^2 + 4}}{2}\nonumber\\
		x_{1,2} = \omega_0^2\frac{\left( 2 - \frac{1}{Q^2} \right)\pm \sqrt{\left( 2 - \frac{1}{Q2} \right)^2 + 4}}{2}
	\end{align}
	Dado que $x=\omega^2$ tenemos que $\omega=\sqrt{x}$, así, tomando la parte positiva de la ecuación anterior obtenemos
	\begin{tcolorbox}[colback=red!5!white,colframe=red!75!black,width=(\linewidth+0.45cm)]
		\begin{align}
			\omega_c = \omega_0\sqrt{\frac{2 - \frac{1}{Q^2} + \sqrt{\left( 2 - \frac{1}{Q^2} \right)^2 + 4} }{2}  }
		\end{align}	
	\end{tcolorbox}

	\subsection{Circuito Sallen Key para el filtro pasa bajas de segundo orden}
	El filtro pasa bajas de segundo orden puede ser diseñado con amplificadores operacionales, 
	como se muestra en la figura 2 este circuito es llamado circuito \textit{Sallen-Key}
	
	Empleando la división de tensión, el voltaje en el nodo $V_c$ es
	\begin{align}
		V_c = \frac{R_B}{R_A + R_B}V_O
	\end{align}
	Donde
	\begin{align}
		K = \frac{R_A + R_B}{R_B} = 1 + \frac{R_A}{R_B}
	\end{align}
	La ecuación (12) puede escribirse como
	\begin{align}
		V_c = \frac{R_B}{R_A + R_B}V_O = \frac{1}{\frac{R_A + R_B}{R_B}} V_O = \frac{V_O}{K}
	\end{align}

	Al aplicar la LCK en el nodo $V_a$ tenemos que
	\begin{align}
		I_{R1} = I_{R_2} + I_{C1}\nonumber
	\end{align}
	\begin{figure}[H]
		\centering
		\includegraphics[scale=1]{figure2}
		\caption{Circuito Sallen-Key para el filtro pasa bajas de segundo orden.}
		\label{fig:pic2}
	\end{figure}
	O bien
	\begin{align}
		\frac{V_{in} - V_a}{R_1} = \frac{V_a - V_b}{R_2} + sC_1(V_a - V_o)\nonumber\\
		R_2V_{in}  R_2V_a = R_1V_a - R_1V_b + sR_1R_2C_1V_a - sR_1R_2C_1V_o
	\end{align}
	
\end{document} 